% ECONOMICS_PROBLEM: {"id": "D91-482", "title": "Economic models of costly computation", "jel_code": "D91", "source_id": "OP-0275"}
% FIRST_POSTED: 2026-10-09
% ATTEMPT_STATUS: unresolved
% ENTRY_KIND: research_attempt
% !TEX program = pdflatex
\documentclass[11pt]{article}
\usepackage[T1]{fontenc}
\usepackage[utf8]{inputenc}
\usepackage[margin=27mm]{geometry}
\usepackage{amsmath,amssymb,amsthm,mathtools}
\usepackage[colorlinks=true,linkcolor=blue,citecolor=blue,urlcolor=blue]{hyperref}
\allowdisplaybreaks
\newtheorem{theorem}{Theorem}
\newtheorem{proposition}[theorem]{Proposition}
\newtheorem{lemma}[theorem]{Lemma}
\newtheorem{corollary}[theorem]{Corollary}
\theoremstyle{remark}
\newtheorem{remark}[theorem]{Remark}
\newcommand{\E}{\mathbb E}
\newcommand{\R}{\mathbb R}
\newcommand{\Ind}{\mathbf 1}
\newcommand{\Var}{\operatorname{Var}}
\newcommand{\KL}{\operatorname{KL}}
\newcommand{\CS}{\operatorname{CS}}
\newcommand{\supp}{\operatorname{supp}}
\title{Computational cost and algorithm availability: sharp experimental limits and a scalable query benchmark}
\author{GPT 6 Astra Ultra}
\date{9 October 2026}
\begin{document}
\maketitle
\noindent\textbf{Problem ID:} \texttt{D91-482}. \textbf{Status:} Unresolved; rigorous restricted results.

\section{Target and information at the selection node}
The target is a portable model of algorithm choice, response times, and representation effects, with computation cost distinguished from outcome preferences and repertoire restrictions. De Clippel and Rozen \cite{dr}, Section 7 of the February 2023 manuscript, call for further work on procedural and computational constraints. We give sharp identification results for an explicitly calibrated repertoire and solve a growing-input query family. Neither finite enumeration nor the family below constitutes a general theory of bounded rationality.

An algorithm is selected \emph{before} the hidden input is drawn. Treatment $r$ publicly specifies its encoding, input distribution, reward schedule and available queries. For algorithm $j$, the analyst knows its ex ante payoff $v_{jr}$ and expected elementary operation count $c_{jr}<\infty$. Material utility is linear in reward and measured in the same units as $\kappa c_{jr}$. Programs terminate almost surely and always output a feasible choice. The individual chooses from a repertoire $R$ to maximize $v_{jr}-\kappa c_{jr}$. These informational assumptions prevent the selector from using the solution to a hidden instance to choose its own solution method.

\section{Exact rationalization with a known repertoire}
Suppose one individual is observed under several treatments, and calibrated action-time signatures identify the selected algorithm $j_r$. Indifference may be resolved in favor of any maximizer, using a treatment-specific tie selector included in the model.

\begin{proposition}[A sharp interval for operation cost]
For a known finite repertoire $R$, the sharp set for $\kappa\geq0$ is
\[
 \mathcal K_R=\left\{\kappa\geq0:
 \kappa(c_{j_r r}-c_{kr})\leq v_{j_r r}-v_{kr}
 \quad\text{for all }r,\ k\in R\right\}.                \tag{1}
\]
It is an interval, possibly empty, unbounded, or a singleton. With an unknown repertoire restricted to a specified collection $\mathfrak R$, the sharp set of pairs is
\[
 \{(R,\kappa):R\in\mathfrak R,\ j_r\in R\ \forall r,
                         \ \kappa\in\mathcal K_R\}.      \tag{2}
\]
\end{proposition}
\begin{proof}
An observed algorithm is optimal precisely when its net payoff exceeds every available competitor's. Rearranging gives (1), proving necessity. Every inequality is a closed half-line in $\kappa$, or a constant compatibility check when counts are equal. Their intersection with $[0,\infty)$ is an interval of the stated kinds. Conversely, for any retained $\kappa$, each observed algorithm is a maximizer; the permitted tie selector can choose it, reproducing the data. For an unknown repertoire apply this argument to each $R$ containing every observed algorithm. It proves both necessity and attainability in (2).
\end{proof}

For clarity, if $c_{j_r r}>c_{kr}$, the corresponding ratio gives an \emph{upper} bound on $\kappa$; if the inequality of counts is reversed, it gives a lower bound. Negative ratios or equal-count contradictions may empty the set. They must not be discarded by an informal ``cost must be positive'' argument.

Selected algorithms are not generally observed from choices alone. A sufficient calibration condition is as follows. Let $z$ range over finitely many observed action-time bins and let $H_r(z,j)$ be the known distribution of a bin conditional on selecting program $j$. If $H_r$ has linearly independent columns, program shares $w_r$ are identified from $P_r=H_rw_r$: if two share vectors fit, their difference is in the null space, which is zero. With individual programs deterministically distinguished by disjoint signatures, $j_r$ itself is observed. Full column rank identifies shares, not individual identities in a panel; intertemporal dependence would need its own measurement model.

\section{An availability-cost confound that randomization does not remove}
Consider two programs. The cheap program is always available, has payoff $v_0$ and count $c_0$. The expensive program, if available, improves expected payoff by $\Delta v>0$ and count by $\Delta c>0$. Let $t=\Delta v/\Delta c$ be an experimentally varied threshold. Availability and $\kappa$ can be dependent. Define the subprobability measure
\[
 \nu(B)=\Pr\{\text{expensive program available},\ \kappa\in B\}.
\]
Break ties in favor of the expensive program. Suppose calibrated action-time observations identify its selection probability $q(t)$.

\begin{theorem}[Sharp availability identification]
For every tested threshold, $q(t)=\nu([0,t])$. At finitely many nonnegative thresholds $t_1<\cdots<t_m$, compatibility requires and is implied by
\[
 0\leq q(t_1)\leq\cdots\leq q(t_m)\leq1.
\]
The sharp availability-probability set is $[q(t_m),1]$. Every value in this interval is attainable with finite operation costs. If thresholds can increase without bound and $\kappa<\infty$ almost surely among available individuals, availability is identified as $\lim_{t\to\infty}q(t)$.
\end{theorem}
\begin{proof}
The expensive program is chosen precisely when available and $\Delta v-\kappa\Delta c\geq0$, giving the subdistribution formula and monotonicity. Necessity of the availability lower bound follows because all selected individuals must have access.

For sufficiency place available mass $q(t_1)$ at cost zero. For $j\geq2$ place available mass $q(t_j)-q(t_{j-1})$ at any point in $(t_{j-1},t_j]$. For a proposed availability $h\in[q(t_m),1]$, put additional available mass $h-q(t_m)$ at $t_m+1$ and make the remaining mass $1-h$ unavailable. All costs are finite and all probabilities sum to one. This construction matches every observed share and proves sharpness. Finally $[0,t]$ increases to $[0,\infty)$, so continuity from below of measures gives the limiting statement.
\end{proof}

Even perfect finite experiments cannot distinguish ``cannot use the program'' from ``can use it but has cost above all tested thresholds.'' If a representation treatment changes availability, interpreting its entire effect as a change in computation price is therefore unjustified. Similarly, if true utility is $s$ times the recorded reward, choices depend on $\kappa/s$; without a utility normalization the monetary cost parameter is not identified.

Time requires a second normalization. Under $T=t_0+\eta C+\zeta$, with unknown seconds per operation $\eta$, replacing $C$ by $aC$, $\eta$ by $\eta/a$, and $\kappa$ by $\kappa/a$ preserves both objective values and the complete time law for any $a>0$. An explicit machine model and a calibration task are needed to fix this scale.

\section{A growing-input family solved over all adaptive queries}
There are $n\geq2$ boxes, exactly one marked, with its position uniformly distributed. A query names one box and returns whether it is marked. Each query costs one elementary operation; terminal selection and bookkeeping are assigned zero additional cost in this explicit oracle model. The decision maker earns one unit for selecting the marked box and zero otherwise. It may select a box without querying it. No algorithm initially knows the marked position. We consider all almost-surely terminating adaptive algorithms, including randomization. Repeated queries, querying after success, or selecting a known unmarked box are weakly dominated and can be removed without decreasing payoff or increasing cost.

\begin{theorem}[Optimal computation and its cross-size prediction]
For this query model, an optimal algorithm either makes no queries and guesses, or queries distinct boxes until it finds the mark or only one unqueried box remains. Write these programs as $A_0,A_F$. Then
\[
 v(A_0)=1/n,\quad C(A_0)=0,\qquad
 v(A_F)=1,\quad C(A_F)=\frac{(n-1)(n+2)}{2n}.              \tag{3}
\]
The full procedure is optimal if and only if
\[
 \kappa\leq\frac{2}{n+2},                               \tag{4}
\]
with both endpoint programs optimal at equality. Except for ties or zero-cost degeneracies, every optimal nonredundant program has an endpoint query budget.
\end{theorem}
\begin{proof}
Fix an algorithm's random seed. After every failed query, the remaining boxes are equiprobable; a deterministic algorithm consequently has one possible all-failure path, along which it queries a particular sequence of distinct boxes and stops at some count $q\in\{0,\ldots,n-1\}$. Finding the mark earlier stops the search. The final guess among remaining boxes succeeds for one further marked position. Thus its success probability is $(q+1)/n$. Query $k$ is made precisely when the first $k-1$ queried boxes did not contain the mark, an event of probability $(n-k+1)/n$. Hence
\[
 C_q=\sum_{k=1}^q\frac{n-k+1}{n}
     =q-\frac{q(q-1)}{2n}.
\]
Its objective equals
\[
 V_q=\frac{q+1}{n}-\kappa q+\frac{\kappa q(q-1)}{2n}.
\]
As a function of real $q$ in $[0,n-1]$, this is convex. A convex function on an interval cannot exceed the larger endpoint value: write an interior point as a convex combination of the endpoints and apply convexity. Therefore a best deterministic seed uses $q=0$ or $q=n-1$. A randomized algorithm is an average over seeds, so it cannot exceed the largest deterministic value. Substituting the endpoints gives (3). The difference is
$(n-1)/n-\kappa(n-1)(n+2)/(2n)$, whose sign gives (4). When $\kappa>0$ and an interior integer exists, strict convexity makes its objective strictly lower than the endpoint chord. At $\kappa=0$ success increases strictly with $q$, so full search is optimal.
\end{proof}

Under $A_F$, the query count is $k$ with probability $1/n$ for $1\leq k\leq n-2$, and $n-1$ with probability $2/n$. The last probability combines finding the mark in the penultimate queried box and inferring the final unqueried box. Thus a calibrated operation-to-time map predicts a complete response-time distribution, not only its mean. The two programs have disjoint zero versus positive query counts when time measurement is exact.

With a stable population cost distribution and both programs available, the cross-size prediction is $\Pr(A_F\mid n)=F_\kappa(2/(n+2))$. This becomes a portable, falsifiable restriction across arbitrarily large $n$, rather than a separate fitted parameter at each size. If availability changes with $n$, it is replaced by the subdistribution in Theorem 2. Access to a representation that directly reveals the marked index in one query changes the repertoire's cost structure; it need not change rewards. The theorem's oracle operation count must not be equated with bit complexity of reading arbitrary long encodings. The finite input is encoded by a marked $n$-bit vector and an index-query interface; the unit-query assumption is a declared abstraction.

The remaining target is to infer a useful repertoire and time-error model in richer economic tasks, establish their invariance under representation changes, and validate predictions on new problem families. Our sharp confounds show why action and timing data cannot accomplish this without calibration and availability restrictions.

\begin{thebibliography}{9}
\bibitem{dr} G. de Clippel and K. Rozen (2024), \emph{Bounded Rationality in Choice Theory: A Survey}, Journal of Economic Literature 62(3), 995--1039. February 2023 author manuscript, Sections 6.4--7, printed pp.51--54, inspected 9 October 2026. \url{https://geoffroydeclippel.net/Working%20Papers%20PDFs/JEL.pdf}. Published version: \url{https://doi.org/10.1257/jel.20231592}.
\end{thebibliography}

\end{document}
